\section{Agentic UI and Web Generation}
\label{sec:ui-web}

Interface generation joins visual appearance to executable behavior. Component trees, style rules, screenshots, browser state, and interaction traces provide complementary views of the same artifact, making this domain a stringent test of whether visual feedback can be translated into code-level repair.

\subsection{From visual specification to executable interface.}
Visual-to-code and specification-to-interface systems must coordinate design intent with component structure and runtime behavior. A rendered screenshot is an observation of the program, not a substitute for functional state.

\paragraph{MM-WebAgent}
MM-WebAgent is a hierarchical multimodal framework that generates coherent webpages with images, videos, and charts from natural language. It operates through global layout planning, local element planning, parallel AIGC-based generation, and three-tier iterative refinement (local asset, surrounding code, and full page) guided by rendered screenshots. ~\cite{avg260415309}.

\paragraph{MAxPrototyper}
MAxPrototyper is a multi-agent system that converts a text description and a wireframe into an editable UI prototype. A Theme Design Agent sets the global theme and generates a theme image via fine-tuned diffusion with ControlNet, while text, image, and icon agents create each component under its direction, using a cache pool to maintain design context.~\cite{yuan2024maxprototyper}.

\paragraph{GameUIAgent}
GameUIAgent is an LLM-powered framework that generates editable Figma game UI designs from natural language via a Design Spec JSON intermediate representation, combining LLM generation, deterministic post-processing, and a VLM-driven Reflection Controller that iteratively refines designs with guaranteed non-regressive quality~\cite{avg260314724}. The Quality Ceiling Effect (r = -0.96) and Rendering-Evaluation Fidelity Principle further reveal that improvement is bounded by evaluator headroom and that incomplete rendering can paradoxically degrade VLM scores.

\subsection{Browser-grounded verification and repair.}
Browser-grounded systems can combine screenshots, DOM structure, interaction tests, and accessibility checks. This heterogeneous evidence makes precise repair possible but also reveals failures that a single image metric cannot capture.


\paragraph{PlayCoder}
PlayCoder is a multi-agent framework that generates executable GUI application code through a closed-loop of repository-aware generation, automated visual testing, and iterative repair, targeting silent behavioral failures that evade compilation and unit tests~\cite{avg260419742}. On the PlayEval benchmark, PlayCoder substantially outperforms baselines, reaching up to 38.1\% Exec@3 and 20.3\% Play@3 with Claude Sonnet 4 and Qwen3-Coder.

\paragraph{TDDev}
 is a test-driven development framework that generates full-stack web applications from natural language and optional design images via three agents: test generation, development, and testing (with BrowserUse for interaction simulation), iteratively refining code until all test cases pass~\cite{avg250925297}.

\paragraph{VISTA}
VISTA is a benchmark evaluating end-to-end web-app coding agents under five input conditions varying visual/structural fidelity and stack constraint, with 3,253 annotated interactions across 128 pages~\cite{avg260526144}. It combines DOM-grounded behavior checks with CLIP-based visual similarity. Among four model–harness systems, free-stack conditions yield the best Combined scores (~0.26). Opus achieves the highest Combined (0.261) and behavior (0.336), while GPT-5.5 attains the best CLIP (0.853) but lowest behavior (0.283), revealing a decoupling between visual reproduction and functional correctness.

\paragraph{Vision2Web}
Vision2Web is a hierarchical benchmark for evaluating multimodal coding agents on visual website development across three levels—static, interactive frontend, and full-stack—with 193 tasks, 918 prototypes, and 1,255 test cases~\cite{avg260326648}. It employs workflow-based verification: a GUI agent checks functionality and a VLM judge assesses visual fidelity. Claude Opus 4.5 performs best overall under OpenHands (full-stack VS/FS: 38.4/57.6), while Gemini-3-Pro-Preview leads on static desktop (63.3) but collapses on full-stack (VS 11.7, FS 22.6). The GUI verifier achieves 87.2\% node-level agreement with human annotations, and the VLM judge attains a mean Spearman correlation of 0.66.

UI and Web work provides the clearest opportunity to connect perception to executable repair because the same artifact can be inspected as pixels, a component tree, and runtime behavior. The principal challenge is to preserve this correspondence across long editing trajectories rather than optimizing isolated screenshots.

\FloatBarrier

\Needspace{0.80\textheight}
